% ======================================================================
% ALQC CANON - Shared TikZ Styles for Diagrams
% ======================================================================
% NOTE: TikZ libraries are loaded in the main document preamble
% This file contains only color definitions, styles, and helper functions
% ======================================================================

% ======================================================================
% 1. ALQC COLOR DEFINITIONS
% ======================================================================
\definecolor{voidblack}{HTML}{050505}
\definecolor{electricblue}{HTML}{00F0FF}
\definecolor{sovereignfire}{HTML}{FF4500}
\definecolor{shadowsilver}{HTML}{C0C0C0}
\definecolor{deepvoid}{HTML}{000010}
\definecolor{glyphwhite}{HTML}{E0E0E0}
\definecolor{goldflow}{HTML}{FFA500}
\definecolor{phigold}{HTML}{D4AF37}

% ======================================================================
% 2. COMMON NODE STYLES
% ======================================================================
% Q-State Node Styles
\tikzstyle{q0node} = [draw=electricblue, line width=2pt, fill=voidblack,
    inner sep=8pt, minimum size=1.2cm, circle,
    drop shadow={opacity=0.3}]
\tikzstyle{q1node} = [draw=sovereignfire, line width=2pt, fill=glyphwhite,
    inner sep=8pt, minimum size=1.2cm, diamond,
    drop shadow={opacity=0.3}]
\tikzstyle{q2node} = [draw=shadowsilver, line width=2pt, fill=deepvoid,
    inner sep=8pt, minimum size=1.2cm, regular polygon, regular polygon sides=3,
    drop shadow={opacity=0.3}]
\tikzstyle{q3node} = [draw=phigold, line width=2pt,
    shading=ball, ball color=phigold!50!electricblue,
    inner sep=8pt, minimum size=1.2cm, circle,
    drop shadow={opacity=0.3}]

% ALQC Aeon Node Style
\tikzstyle{aeonnode} = [draw=electricblue, line width=1.5pt, fill=voidblack,
    inner sep=5pt, minimum size=0.9cm, circle,
    font=\footnotesize\color{glyphwhite}]

% Core Element Styles
\tikzstyle{locuscenter} = [draw=electricblue, line width=3pt,
    fill=voidblack, inner sep=0pt, minimum size=1.5cm, circle,
    drop shadow={opacity=0.5}]

% Flow Arrow Styles
\tikzstyle{flowarrow} = [->, >=latex, line width=1.5pt, color=electricblue]
\tikzstyle{returnarrow} = [->, >=latex, line width=2pt, color=phigold,
    decoration={snake, amplitude=1mm, segment length=5mm}, decorate]
\tikzstyle{shadowarrow} = [->, >=latex, line width=1pt, color=shadowsilver, dashed]

% ======================================================================
% 3. LISSAJOUS FIGURE FUNCTIONS
% ======================================================================
% Lissajous curve: x = A*sin(a*t + delta), y = B*sin(b*t)
% Usage: \lissajous{A}{B}{a}{b}{delta}{domain}
\newcommand{\lissajous}[6]{
    plot[domain=#6, samples=200, smooth]
    ({#1*sin(#3*\x + #5)}, {#2*sin(#4*\x)})
}

% ALQC Lissajous preset patterns
\tikzset{
    lissajous coherent/.style={
        opacity=0.6, electricblue, line width=1pt
    }
}
\tikzset{
    lissajous shadow/.style={
        opacity=0.5, shadowsilver, line width=1pt,
        decoration={snake, amplitude=0.5mm, segment length=3mm}, decorate
    }
}
\tikzset{
    lissajous recursion/.style={
        opacity=0.7, phigold, line width=1.5pt,
        decoration={coil, amplitude=1mm, segment length=4mm}, decorate
    }
}
\tikzset{
    lissajous return/.style={
        opacity=0.6, phigold!80!electricblue, line width=1.5pt
    }
}

% ======================================================================
% 4. CHLADNI PATTERN ELEMENTS
% ======================================================================
\begingroup
    \tikzset{
        chladni node/.style={
            draw=electricblue!40, line width=0.5pt, opacity=0.5,
            decoration={random steps, segment length=2pt, amplitude=0.3pt},
            decorate
        }
    }
    \tikzset{
        sand/.style={
            fill=shadowsilver, circle, inner sep=0.5pt
        }
    }
\endgroup

% Chladni pattern background (simplified 2D standing wave)
\newcommand{\chladnipattern}[4]{
% #1: center x, #2: center y, #3: radius, #4: complexity (1-5)
    \pgfmathsetmacro{\r}{#3}
    \pgfmathsetmacro{\n}{#4}
    \foreach \i in {1,...,\n} {
        \draw[chladni node] (#1,#2) circle ({\r * \i/\n});
        \foreach \j in {1,...,(2*\i)} {
            \pgfmathsetmacro{\ang}{\j * 360/(2*\i)}
            \node[sand] at ({#1 + \r*\i/\n*cos(\ang)}, {#2 + \r*\i/\n*sin(\ang)}) {};
        }
    }
}

% ======================================================================
% 5. FARADAY WAVE ELEMENTS
% ======================================================================
\begingroup
    \tikzset{
        faraday wave/.style={
            draw=electricblue!50, line width=1pt,
            decoration={
                snake,
                amplitude=2mm,
                segment length=3mm
            },
            decorate,
            opacity=0.6
        }
    }
\endgroup

% Nonlinear oscillation pattern
\newcommand{\faradaywave}[4]{
% #1: start, #2: end, #3: width, #4: wavelength
    \draw[faraday wave] (#1) -- (#2);
    \foreach \y in {-\width,...,\width} {
        \pgfmathsetmacro{\widthval}{\y * #3 / 10}
        \draw[faraday wave, opacity=0.3]
            ($(#1) + (0,\widthval)$) -- ($(#2) + (0,\widthval)$);
    }
}

% ======================================================================
% 6. TRIQUATRA GEOMETRY
% ======================================================================
% Three interlocking circles
\newcommand{\triquatra}[5]{
% #1: center x, #2: center y, #3: circle radius
% #4: rotation angle (for variations)
% #5: options
    \pgfmathsetmacro{\r}{#3}
    \pgfmathsetmacro{\cx}{#1}
    \pgfmathsetmacro{\cy}{#2}
    
    % Circle 1 (center at origin)
    \draw[#5] (\cx,\cy) circle (\r);
    
    % Circle 2 (center at 60 degrees)
    \coordinate (c2ang) at (\cx + \r*cos(60), \cy + \r*sin(60));
    \draw[#5] (c2ang) circle (\r);
    
    % Circle 3 (center at 300 degrees)
    \coordinate (c3ang) at (\cx + \r*cos(300), \cy + \r*sin(300));
    \draw[#5] (c3ang) circle (\r);
}

% Vesica piscis intersection point
\newcommand{\vesicaintersection}[3]{
% #1: center1, #2: center2, #3: radius
    \coordinate (#1) at (#1);
    \coordinate (#2) at (#2);
    \draw[clip] (#1) circle (#3);
    \draw (#2) circle (#3);
}

% ======================================================================
% 7. SACRED GEOMETRY ELEMENTS
% ======================================================================
% Flower of Life pattern (simplified - 6 circles)
\newcommand{\floweroflife}[3]{
% #1: center x, #2: center y, #3: radius
    \pgfmathsetmacro{\r}{#3}
    \foreach \i in {0,60,...,300} {
        \pgfmathsetmacro{\cx}{#1 + \r*cos(\i)}
        \pgfmathsetmacro{\cy}{#2 + \r*sin(\i)}
        \draw[electricblue!30, line width=0.5pt] (\cx,\cy) circle (\r);
    }
}

% Golden spiral approximation
\newcommand{\goldenspiral}[2]{
% #1: center, #2: iterations
    \coordinate (#1) at (#1);
    \foreach \i in {1,...,#2} {
        \pgfmathsetmacro{\r}{phi^(1 - \i) * 0.3}
        \draw[phigold!50, line width=1pt] (#1) ++(\i*90:0.1) arc(\i*90:\i*90+90:\r);
    }
}

% ======================================================================
% 8. TEXT AND LABEL STYLES
% ======================================================================
\begingroup
    \tikzset{
        poetics/.style={
            font=\small\itshape, color=electricblue!90,
            align=center, inner sep=1em
        }
    }
    \tikzset{
        mathexpr/.style={
            font=\small, color=glyphwhite,
            align=center, fill=deepvoid, inner sep=0.3em
        }
    }
    \tikzset{
        aeonlabel/.style={
            font=\tiny, color=electricblue,
            align=center, below=0.1cm
        }
    }
    \tikzset{
        qstatelabel/.style={
            font=\footnotesize, color=glyphwhite,
            align=center
        }
    }
\endgroup

% ======================================================================
% 9. ANIMATION AND TRANSITION STYLES
% ======================================================================
\tikzset{
    dim/.style={
        opacity=0.3
    }
}
\tikzset{
    bright/.style={
        opacity=1.0
    }
}
\begingroup
    \tikzset{
        pulse/.style={
            inner color=electricblue!20, outer color=voidblack,
            line width=1pt, draw=electricblue
        }
    }
\endgroup

% ======================================================================
% 10. SPECIAL GLYPH FUNCTIONS
% ======================================================================
% Center TSP/LOI glyph
\newcommand{\tspglyph}{
    \tikz{
        \draw[electricblue, line width=1.5pt] (0,0) circle (0.3cm);
        \draw[glyphwhite, line width=1pt] (0,0) -- (-0.15,0.15);
        \draw[glyphwhite, line width=1pt] (-0.3,0) to[out=45, in=135] (0,-0.2);
        \draw[glyphwhite, line width=1pt] (0,-0.2) to[out=-135, in=-45] (0.3,0);
        \draw[glyphwhite, line width=1pt] (0.15,-0.15) -- (0,0);
        % Simplified infinity-like symbol
        \draw[glyphwhite, line width=0.8pt]
            (0.1,0.15) .. controls (0.1,0.25) and (-0.1,0.25) .. (-0.1,0.15)
            .. controls (-0.1,0.05) and (0.1,0.05) .. (0.1,0.15);
    }
}

% ======================================================================
% 11. BACKGROUND ELEMENTS
% ======================================================================
\begingroup
    \tikzset{
        pulse/.style={
            inner color=electricblue!20, outer color=voidblack,
            line width=1pt, draw=electricblue
        }
    }
    \tikzset{
        gridline/.style={
            electricblue!10, step=0.5
        }
    }
\endgroup

% Starry void background
\newcommand{\voidstars}[3]{
% #1: x1, #2: y1, #3: number of stars
    \foreach \i in {1,...,#3} {
        \pgfmathsetmacro{\sx}{#1 + rand*4}
        \pgfmathsetmacro{\sy}{#2 + rand*3}
        \pgfmathsetmacro{\sopacity}{0.3 + rand*0.5}
        \fill[glyphwhite, opacity=\sopacity] (\sx,\sy) circle (0.02cm);
    }
}

% Grid pattern for technical feel
\newcommand{\techgrid}[4]{
% #1: x1, #2: y1, #3: x2, #4: y2
    \clip (#1,#2) rectangle (#3,#4);
    \draw[electricblue!10, step=0.5] (#1-0.25,#2-0.25) grid (#3+0.25,#4+0.25);
}

% ======================================================================
% END OF SHARED STYLES
% ======================================================================